% from ega4-8-fr.tex
\oldpage[IV-3]{46}
\begin{proof}
En vertu de (8.12.4, (ii)) et du caractère local (sur $Y$) des notions de
morphisme quasi-fini, séparé et de présentation finie, on peut se borner au
cas où $Y=\operatorname{Spec}(A)$ est affine. Appliquant (8.9.1), on peut
supposer qu'il y a un sous-anneau $A_0$ de $A$, qui est une
$\mathbf Z$-algèbre de type fini, et un $A$-isomorphisme
$X_0\otimes_{A_0}A\xrightarrow{\sim}X$, $f$ s'identifiant par cet
isomorphisme à $f_0\times 1$, où
$f_0:X_0\to\operatorname{Spec}(A_0)$ est un morphisme \emph{de type fini}~;
en outre (8.10.5, (v) et (xii)) on peut supposer que $f_0$ est séparé et
quasi-fini~; si l'on prouve que $f_0$ est pseudo-fini, il en sera de même de
$f$ d'après (8.12.4, (i)). Comme $A_0$ est alors noethérien et que les
notions de morphisme de type fini, séparé et quasi-fini se conservent par
changement de base, il résulte de (8.12.4, (iii)) que l'on peut même supposer
que $A$ est un anneau local, \emph{essentiellement de type fini sur}
$\mathbf Z$ (1.3.8). Posons $n=\dim(A)$, et procédons par récurrence sur $n$~;
pour $n=0$, le théorème est évident, $A$ étant un corps et le morphisme $f$
étant déjà fini (II, 6.2.2). Posons $B=\Gamma(X,\mathcal O_X)$~; désignons
par $C$ la fermeture intégrale de $A$ dans $B$, posons
$Z=\operatorname{Spec}(C)$ et soit $g:X\to Z$ le $Y$-morphisme correspondant
au $A$-homomorphisme canonique injectif $C\to B$~; en vertu de (8.12.3), il
s'agit de montrer que $g$ est une immersion ouverte. Soit $a$ le point fermé
de $Y$, et soit $U=Y-\{a\}$~; $U$ est noethérien et tous ses anneaux locaux
sont essentiellement de type fini sur $\mathbf Z$ et de dimension $<n$~;
compte tenu de l'hypothèse de récurrence, et de (8.12.4, (iii)), on voit que
la restriction $f^{-1}(U)\to U$ de $f$ est un morphisme pseudo-fini. On en
conclut (8.12.3) que, si $h:Z\to Y$ est le morphisme structural, la
restriction $f^{-1}(U)\to h^{-1}(U)$ de $g$ est une immersion ouverte.
Posons $A'=\widehat A$, $Y'=\operatorname{Spec}(A')$,
$X'=X_{(Y')}$, $f'=f_{(Y')}:X'\to Y'$. Comme le morphisme canonique
$u:Y'\to Y$ est plat, il résulte de (2.3.1) que
$B'=\Gamma(X',\mathcal O_{X'})$ s'identifie à la $A'$-algèbre
$B\otimes_A A'$.

D'autre part, comme $A$ est un anneau local \emph{excellent} (7.8.3), le
morphisme $u:Y'\to Y$ est régulier, et \emph{a fortiori} normal, et par suite
(6.14.4) \emph{la fermeture intégrale de $A'$ dans $B'$ est égale à}
$C'=C\otimes_A A'$. On voit donc que $Z'=\operatorname{Spec}(C')$ est égal à
$Z_{(Y')}$ et le morphisme $g':X'\to Z'$ provenant de l'injection $C'\to B'$
est égal à $g_{(Y')}$. Comme $u:Y'\to Y$ est fidèlement plat et
quasi-compact, pour prouver que $g$ est une immersion ouverte, il suffit de
prouver que $g'$ est une immersion ouverte (2.7.1, (x)). Notons maintenant
que $u^{-1}(a)$ est réduit au point fermé $a'$ de $Y'$ et par suite
$U'=Y'-\{a'\}=u^{-1}(U)$. Si $h':Z'\to Y'$ est le morphisme canonique, le
fait que la restriction $f^{-1}(U)\to h^{-1}(U)$ de $g$ soit une immersion
ouverte entraîne qu'il en est de même de la restriction
$f'^{-1}(U')\to h'^{-1}(U')$ de $g'$. Notons maintenant que $f'$ est un
morphisme séparé et quasi-fini (II, 6.2.4)~; comme $A'$ est \emph{complet},
on déduit de (II, 6.2.6) que $X'$ est $Y'$-isomorphe à une somme
$X'_1\amalg X'_2$, où la restriction
$f'|X'_1=f'_1:X'_1\to Y'$ est un morphisme \emph{fini}, et
$X'_2\subset f'^{-1}(U')$. Il en résulte que $B'$ est composée directe des
deux $A'$-algèbres
$\Gamma(X'_1,\mathcal O_{X'_1})=B'_1$ et
$\Gamma(X'_2,\mathcal O_{X'_2})=B'_2$~; on en conclut aussitôt que la
fermeture intégrale $C'$ de $A'$ dans $B'$ est composée directe des
fermetures intégrales $C'_1$, $C'_2$ de $A'$ dans $B'_1$, $B'_2$
respectivement, d'où $Z'=Z'_1\amalg Z'_2$, où
$Z'_i=\operatorname{Spec}(C'_i)$ ($i=1,2$)~; et le morphisme canonique
$g':X'\to Z'$ est tel que $g'|X'_i$ soit le morphisme canonique
$g'_i:X'_i\to Z'_i$ ($i=1,2$). Mais comme $B'_1$ est déjà une
$A'$-algèbre finie, on a $C'_1=B'_1$, et $g'_1$ est donc un isomorphisme.
D'autre part, comme $X'_2\subset f'^{-1}(U')$ et est ouvert dans
$f'^{-1}(U')$, on sait
